\documentclass[11pt]{article}
\usepackage[margin=0.9in]{geometry}
\usepackage{amsmath,amssymb}
\usepackage[T1]{fontenc}
\usepackage{lmodern}
\usepackage{xcolor}
\usepackage{longtable}
\usepackage{enumitem}

\definecolor{errrule}{RGB}{90,90,140}
\newcommand{\ket}[1]{\left|{#1}\right\rangle}
\newcommand{\bra}[1]{\left\langle{#1}\right|}
\newcommand{\LR}{\ket{L}\!\bra{L}}

\setlength{\parindent}{0pt}

\begin{document}

\begin{center}
{\LARGE\bfseries Errata sheet --- corrected unified edition}\\[6pt]
{\large A Theory of Everything Research Program (Baird \& ZoraASI, 2026)}\\[2pt]
{\large Inserted page \#1; corpus pages 1--6894 unchanged; pages 6895--6909 replaced by the corrected\\[-2pt] methods note of 11 September 2026.}\\[4pt]
{\small Produced 11 September 2026. This sheet documents every change made in the corrected edition.}
\end{center}

\vspace{4pt}
\rule{\linewidth}{0.5pt}

\section*{A list of every change}

Each entry states the original printed expression, the corrected expression, and the reason. The
corrected note itself flags each repair with a boxed \textbf{Correction note}; this sheet is the
stand-alone index for the boundary between corpus pages 6894 and the corrected note.

\begin{longtable}{p{0.9in}p{2.35in}p{2.35in}p{2.35in}}
\hline
\textbf{Location} & \textbf{Original} & \textbf{Corrected} & \textbf{Why}\\
\hline
\endfirsthead
\hline
\textbf{Location} & \textbf{Original} & \textbf{Corrected} & \textbf{Why}\\
\hline
\endhead

Eq.~(11) alone &
\hspace*{-3pt}$H_{SR} = \lambda(E)\,\LR \otimes B$; single monitor on path $L$ &
\hspace*{-3pt}$H_{SR} = \lambda(\Delta E)\,\LR \otimes B$, $\Delta E \equiv E(x_L)-E(x_R)$; twin monitors on both arms &
Setting $E_L=E_R$ does not remove the printed coupling. Differential gate makes N1 a genuine null.\\

Table~2, row N1 &
$\alpha = 0$ ``for a path-local gate'' &
$\alpha = 0$ because $\lambda(0)=1$ at $\Delta E=0$ (monitor-on/$E{=}0$ reference) &
Null is now a theorem of the model, not an unsupported assertion.\\

Eq.~(18) and (19) &
(18) $2\pi\lim_{\omega\to0^+} J(\omega)\coth(\beta\omega/2)$; (19) $\Gamma_\phi^{(0)} = 2\pi\,\eta_{\mathrm{bath}}k_BT$ &
(19) $\Gamma_\phi^{(0)} = 4\pi\,\eta_{\mathrm{bath}}k_BT$; coherence rate $\tfrac12\Gamma_\phi^{(0)}=2\pi\,\eta_{\mathrm{bath}}k_BT$ &
Printed Ohmic $J(\omega){=}\eta\omega\,e^{-\omega/\omega_c}$ into (18) gives $4\pi\eta k_BT$, not $2\pi$. Spectral convention reconciled.\\

Eq.~(21)/(22) and \S9 baseline &
\eqref{eq:contrast} ``$V/V_{\mathrm{QM}}=\exp[-\tfrac12\Gamma_\phi^{(0)}T(\lambda^2-1)]$'' with ``if desired'' $E{=}0$ clause; \S9 used N0 (monitor-off) &
Two named references: $V_{\mathrm{QM}}^{\mathrm{(on)}}$ (monitor-on, $E{=}0$) and $V_{\mathrm{QM}}^{\mathrm{(off)}}$ (monitor-off). The ``$-1$'' is bound to the monitor-on reference; estimator declares the baseline &
The $-1$ requires a monitor-on/$\lambda{=}1$ reference; it does not apply unchanged to a monitor-off baseline.\\

\S9.2 estimator &
$\hat V_{\mathrm{QM}}$ from ``contemporaneous N0 (or N1 after N1 is shown dead)'' &
Baseline declared per row: default monitor-on (N1); N0 is the monitor-off offset meter &
Prevents the printed conflation of two different comparisons.\\

\S9.4 kill-plot, rule 2 &
``N1 sits on the same band as N0'' &
``N1 sits on 0 (monitor-on, $\Delta E{=}0$ null)''; N0 sits on its monitor-off band &
N1 vs N0 are different baselines; comparing them directly could false-fail the mechanism.\\

\S9.5 analyzer predicate &
``$|\hat\alpha_{N1}-\hat\alpha_{N0}| > \alpha_{\mathrm{CM}}^{\max}$'' &
Absolute N1 null band plus a recorded baseline column &
Compares two different baselines in the original; now unambiguous.\\

Appendix A sentence &
``$\Gamma_\phi^{(0)} = \lambda(\Delta E)^2\, 2\pi\,\eta_{\mathrm{bath}}k_BT$'' &
``$\Gamma_\phi^{(0)} = \lambda(\Delta E)^2\, 4\pi\,\eta_{\mathrm{bath}}k_BT$'' &
Consistency with corrected eq.~(19).\\
\hline
\end{longtable}

\section*{Page-map of the corrected edition}

\begin{center}
\begin{tabular}{lcl}
\textbf{Page(s)} & \textbf{Content} & \textbf{Status}\\
\hline
\#1 & This errata sheet & New\\
1--6894 & Corpus pages (incl. Sept 2026 status note, Part 0) & Unchanged\\
6895--6909 & Reservoir/protocol methods note, 11 Sep 2026 & \textbf{Replaced} by corrected note (14 pp.)\\
\hline
\end{tabular}
\end{center}

\section*{Notes}

\begin{itemize}
\itemsep1pt
\item Historical pages 1--6894 are passed through byte-identical by page selection (qpdf); none are re-typeset.
\item The corrected note renumbers nothing---its equation numbers match the printed note's numbers (11, 14, 18, 19, 20, 21, 22, 23--26, 29--34, both twin models, Table 2) so old cross-references to ``eq.~21'', ``Table 2'', ``Appendix A'' still resolve.
\item One presentation edit of the ledger table (``Still required'' column) corrects a truncated extraction artifact in the compiled corpus; no physics is changed there.
\item No other page in the corpus contains the three corrected equations (verified by full-corpus text audit, 6{,}909 pages); remote references to the note are therefore consistent with the corrected edition.
\end{itemize}

\end{document}